\section{Weekly Summary}

\begin{tcolorbox}[colback=yellow!10,colframe=black!50,title=AI-Generated Content]
\noindent The summary below was written by Claude (Anthropic) from that week's regression outputs and series descriptions. It has not been reviewed or fact-checked by a human, and should be read as an automated, informal recap -- not analysis.
\end{tcolorbox}

Welcome back to the weekly ritual where we take two economic time series that have never met, forcibly introduce them, and then act shocked when they occasionally hit it off. As always, the house motto applies: these pairs are drawn at random, nobody is claiming anything causes anything, and if you find yourself nodding along to a conclusion, please lie down until the feeling passes.

We opened Monday with a genuine plot twist: US Acquisition of Portfolio Investment Assets versus some long-discontinued Freddie Mac mortgage pool series, and for once the universe gave us nothing. The raw regression limped in with an R-squared of basically one percent and a p-value north of 0.3, meaning even before detrending we had a relationship with the structural integrity of a wet paper towel. The detrended version somehow found a way to do worse. A rare clean miss on both counts -- honestly refreshing, like a horoscope that just says "nothing will happen today." Then Tuesday handed us the week's great cautionary tale: Idaho's minimum wage versus Exports of Other Private Services, with a raw R-squared of 0.87 and a p-value so small it needs a telescope. Breathtaking! Idaho's hourly wage apparently explains the nation's service exports! Except, of course, both numbers just drift upward over the decades like everything else does, and once we strip out that shared march of time the R-squared collapses to 0.28 -- technically still "significant," but a shadow of its former swagger. Classic trend mirage, partially debunked.

Wednesday's pairing of a Transportation Services CPI against restaurant and bar sales looked even more spectacular in the raw -- an R-squared of 0.95, the kind of number that makes you want to call your old statistics professor and apologize. But again, two things that both go up over thirty years will cuddle for purely chronological reasons. Detrending knocked it down to 0.47, which, I'll grudgingly admit, is one of the more durable survivors of the week; both transport costs and dining out do respond to the broader cost of living, so a lingering link isn't the most insane thing ever uttered. I remain legally obligated to remind you this is a coincidence factory. Thursday's "arts and recreation job openings versus Real MZM Money Stock" was pure comedy: a shiny raw R-squared of 0.26 that detrending vaporized into a p-value of 0.71, the statistical equivalent of a magician whose rabbit was just a shared upward trend the whole time.

The real standout, if we must crown one, is Friday's Waterloo-Cedar Falls, Iowa house price index against the NAFTA composite leading indicator. The raw R-squared was a jaw-dropping 0.97 -- I briefly considered drafting a Nobel acceptance speech -- and even after detrending it held at a respectable 0.31 with a p-value that refused to die. That's the rarer, more interesting kind of result: a relationship that survives having its trend surgically removed, probably because regional home prices and broad economic momentum genuinely share some underlying business-cycle rhythm. Or it's noise wearing a convincing costume; with this project you truly never know. The week wound down gently with Mideast amusement-and-gambling GDP versus South African construction employment surveys (nothing, raw or detrended, a merciful dud) and Real Consumption of Fixed Capital versus a discontinued Korean employment survey, which squeaked past the raw significance line at p=0.045 before politely fading to 0.059 under detrending -- the statistical version of someone almost getting away with it. All in all: two findings with a pulse, several beautiful mirages, and absolutely no reason to rearrange your portfolio around Idaho's minimum wage. See you next week.
